\documentclass{beamer}

\usepackage{../../../latex_style/beamerthemeExecushares}
\usepackage{../../../latex_style/notations}

\title{Lecture 8.1: Singular Value Decomposition}
\subtitle{Optimization and Computational Linear Algebra for Data Science}
\author{Léo Miolane}
\date{}

\setcounter{showSlideNumbers}{1}

\begin{document}
\setcounter{showProgressBar}{0}
\setcounter{showSlideNumbers}{0}

\frame{\titlepage}
\setcounter{framenumber}{0}
%\setcounter{showProgressBar}{1}
\setcounter{showSlideNumbers}{1}



\begin{frame}[t]{PCA}
	\grid
	
	\begin{itemize}
		\item Data matrix \quad $A \in \R^{n \times m}$
			\vspace{1mm}
		\item ``Covariance matrix''\quad $S = A^{\sT} A \in \R^{m \times m}$.
			\vspace{1mm}
		\item $S$ is symmetric positive semi-definite.
			\vspace{1mm}
		\item \textbf{Spectral Theorem:} there exists an orthonormal basis $v_1, \dots, v_m$ of $\R^m$ such that the $v_i$'s are eigenvectors of $S$ associated with the eigenvalues $\lambda_1 \geq \cdots \geq \lambda_m \geq 0$.
	\end{itemize}

\end{frame}

\begin{frame}[t]{Singular values/vectors}
	\grid

\end{frame}
\begin{frame}[t]{Singular Value decomposition}
	\grid

	\vspace{-0.3cm}
	\begin{block}{Theorem}
		Let $A \in \R^{n \times m}$. Then there exists two orthogonal matrices $U \in \R^{n \times n}$ and $V \in \R^{m \times m}$ and a matrix $\Sigma \in \R^{n \times m}$ such that $\Sigma_{1,1} \geq \Sigma_{2,2}  \geq \cdots \geq 0$ and $\Sigma_{i,j} = 0$ for $i\neq j$, that verify
		$$
		A = U \Sigma V^{\sT}.
		$$
	\end{block}

\end{frame}

\begin{frame}[t]{Geometric interpretation of \ $U\Sigma V^{\sT}$}
\begin{figure}[H]
	\begin{center}
	\includegraphics[width = 0.6\linewidth]{./svd.pdf}
	\end{center}
\end{figure}

\end{frame}

\begin{frame}[t]{Example}
	\grid

\end{frame}


\end{document}
